%%%%% Základní nastavení pro jednostranný tisk:
%%%%% ----------------------------------------------------
% Okraje: levý 40mm, pravý 25mm, horní a dolní 25mm (ale pozor, LaTeX si sám přidává 1in)
\documentclass[12pt, a4paper]{report}
\usepackage{geometry}
\setlength\textwidth{145mm}
\setlength\textheight{247mm}
\setlength\oddsidemargin{15mm}
\setlength\evensidemargin{15mm}
\setlength\topmargin{0mm}
\setlength\headsep{0mm}
\setlength\headheight{0mm}
% \openright zařídí, aby následující text začínal na pravé straně knihy
\newcommand{\openright}{\clearpage}


%%%%% Základní nastavení pro oboustranný tisk:
%%%%% ----------------------------------------------------
% \documentclass[12pt, a4paper, twoside, openright]{report}
% \setlength\textwidth{145mm}
% \setlength\textheight{247mm}
% \setlength\oddsidemargin{15mm}
% \setlength\evensidemargin{0mm}
% \setlength\topmargin{0mm}
% \setlength\headsep{0mm}
% \setlength\headheight{0mm}
% \let\openright=\cleardoublepage


%%%%% Nastavení kódování vstupních souborů: UTF-8
%%%%% ---------------------------------------------------------------
\usepackage[utf8]{inputenc}


%%%%% Nastavení češtiny (slovenština analogicky)
%%%%% ---------------------------------------------------------------

%%% Existují dvě hlavní možnosti, jak zacházet s češtinou. Je zapotřebí zvolit právě jednu.
%%%

%%% MOŽNOST 1 (doporučujeme):
%%% * použití balíčku czech
%%%   (mimo jiné již obsahuje příkaz \uv pro sazbu českých uvozovek)
%%% * kompilace musí následně probíhat pomocí CSLaTeXu (příkaz
%%%   cslatex, resp. cspdflatex)
% \usepackage{czech}

%%% MOŽNOST 2: (zde zakomentovaná)
%%% * použití balíčku babel s volbou pro češtinu
%%% * kompilace následně probíhá standardním LaTeXem (příkaz latex,
%%% resp. pdflatex)
\usepackage[czech]{babel}
\ifx\uv\undefined\newcommand{\uv}[1]{,,#1``}\fi
%%% příkaz pro sazbu českých/slovenských uvozovek
%%% (v novějších verzích babelu je již k dispozici, stejně tak je již
%%% k dispozici v balíčku czech)


%%% Další užitečné balíčky (jsou součástí běžných distribucí LaTeXu)
%%% ----------------------------------------------------------------
\usepackage{amsmath}                %%% rozšíření pro sazbu matematiky
\usepackage{amsfonts}               %%% matematické fonty
\usepackage{amsthm}                 %%% sazba vět, definic apod.
\usepackage{bm}                     %%% tučné symboly (příkaz \bm)
\usepackage{graphicx}               %%% vkládání obrázků
\usepackage{psfrag}                 %%% dodatečná úprava popisků v postscriptových obrázcích
\usepackage{fancyvrb}               %%% vylepšené prostředí pro strojové písmo
\usepackage{natbib}                 %%% zajištuje možnost odkazovat na reference stylem AUTOR (ROK), resp. AUTOR [ČÍSLO]
\usepackage{tikz}                   %%% vkládání vektorových obrázků
\usepackage{bbding}                 %%% balíček s nejrůznějšími symboly
\usepackage{icomma}                 %%% inteligetní čárka v matematickém módu
\usepackage{dcolumn}                %%% lepší zarovnání sloupců v tabulkách
\usepackage{booktabs}               %%% lepší vodorovné linky v tabulkách
\usepackage{paralist}               %%% lepší enumerate a itemize
\usepackage{indentfirst}            %%% zaveď odsazení 1. odstavce
\usepackage[nottoc]{tocbibind}      %%% zajistí přidání seznamu literatury, obrázků a tabulek do obsahu
\usepackage[unicode]{hyperref}      %%% zajištuje generování hyperodkazů, bookmarků atp.
\hypersetup{pdftitle=Webové řešení na prodej vstupenek s~rezervací míst,
    pdfauthor=Filip Ditrich
    ps2pdf,
    colorlinks=false,               %% hyperlinky budou označeny červenými rámečky, které budou neviditelné při tisku na papír
    urlcolor=blue,
    pdfstartview=FitH,
    pdfpagemode=UseOutlines,
    pdfnewwindow,
    breaklinks                      %% zajistí, aby se dlouhé hyperodkazy mohly lámat přes více řádků
}
\usepackage{titlesec}
\titlespacing*{\chapter}{0pt}{-10mm}{5mm}
\titleformat{\chapter}{\normalfont\huge\bfseries}{\thechapter}{1em}{}


%%% Příkazy zjednodušující přenositelnost
%%% -------------------------------------
\newcommand{\FIGDIR}{./figures}    %%% cesta do adresare s obrazky


%%% Seznam použité literatury
%%% Příkaz \bibliographystyle určuje, jakým stylem budou citovány
%%% odkazy v textu, a podle jakého stylu se automaticky vygeneruje
%%% seznam literatury. V závorce je název zvoleného .bst souboru.
%%% Styly plainnat a unsrt jsou standardní součástí latexových
%%% distribucí. Styl czplainnat vyžaduje přítomnost souboru
%%% czplainnat.bst ve stejném direktoráři, v němž se nachází
%%% kompilovaná práce.
%%%
%%% Seznam literatury se vytváří na konci práce příkazem \bibliography, kde v závorce
%%% uvádíme název databázového bib souboru.
%%%
%%%
\bibliographystyle{czplainnat}    %% Autor (rok) s českými spojkami
%\bibliographystyle{plainnat}     %% Autor (rok) s anglickými spojkami
%\bibliographystyle{unsrt}        %% [číslo]

\renewcommand{\bibname}{Seznam použité literatury}


%%%%% Použití fancyvrb (fancy verbatim) při definici prostředí pro
%%%%% sazbu kódu, resp. výstupů z počítačových programů
%%%%% ------------------------------------------------------------
\DefineVerbatimEnvironment{PCinout}{Verbatim}{fontsize=\small, frame=single}


%%%%% Hlavní část dokumentu
%%%%% ---------------------
\begin{document}
%%% titulní strana
    %%%
%%%  TITULNÍ STRANA
%%%
%%%  * soubor obsahující titulní stranu bakalářské práce
%%%
%%%  ===========================================================================
\pagestyle{empty}
\begin{center}

%%% název školy
{\bfseries\large UNICORN VYSOKÁ ŠKOLA s.r.o.}

    \vspace{5mm}

    %%% název oboru
    {\Large Softwarový vývoj}

    \vfill
    \vspace{5mm}

    %%% logo školy
    \centerline{\mbox{\includegraphics[width=83.3mm]{\FIGDIR/uu-icon}}}

    \vfill
    \vspace{5mm}

    %%% typ práce
    {\large\MakeUppercase{Bakalářská práce}}

    \vspace{15mm}

    %%% název práce
    {\LARGE\bfseries Webové řešení na prodej vstupenek s~rezervací míst}

    \vfill

    %%% autor a vedoucí práce
    \begin{tabular}{rl}
        Autor bakalářské práce:   & Filip Ditrich              \\
        \noalign{\vspace{2mm}}
        Vedoucí bakalářské práce: & Ing.\ Marek Beránek, Ph.D. \\
    \end{tabular}

    \vfill

    %%% rok
    Praha 2023

\end{center}

%%% TODO: kopie zadání
%%% čestné prohlášení
    %%%
%%%  ČESTNÉ PROHLÁŠENÍ
%%%
%%%  * soubor obsahující čestné prohlášení k bakalářské práci
%%%
%%%  ===========================================================================
\newpage
\pagestyle{empty}
\vspace*{\stretch{8}}

%%% nadpis
\noindent
{\large\bfseries Čestné prohlášení}\\

%%% text
\noindent
Prohlašuji, že jsem svou bakalářskou práci na~téma \textit{Webové řešení na prodej vstupenek s~rezervací míst} vypracoval samostatně pod~vedením vedoucího bakalářské práce a s~použitím výhradně odborné literatury a~dalších informačních zdrojů, které jsou v práci všechny citovány a~jsou také uvedeny v~seznamu použitých zdrojů.\\

\noindent
Jako autor této bakalářské práce dále prohlašuji, že v~souvislosti s~jejím vytvořením jsem neporušil autorská práva třetích osob a~jsem si plně vědom následků porušení ustanovení § 11 a následujících autorského zákona č.~121/2000~Sb.\\

\noindent
Dále prohlašuji, že odevzdaná tištěná verze bakalářské práce je shodná s~verzí, která byla odevzdána elektronicky.

%%% podpis - místo/den
\vspace{18mm}
\noindent
V \makebox[4cm]{\dotfill} dne \makebox[2.5cm]{\dotfill}
\hspace*{\fill}
\makebox[4cm]{\dotfill}

%%% podpis
\begin{flushright}
    \noindent
    Filip Ditrich
\end{flushright}

%%% poděkování
    %%%
%%%  PODĚKOVÁNÍ
%%%
%%%  * soubor obsahující poděkování za pomoc při vytvoření bakalářské práce
%%%
%%%  ===========================================================================
\newpage
\pagestyle{empty}
\vspace*{\stretch{8}}

%%% nadpis
\noindent
{\large\bfseries Poděkování}\\

%%% TODO: text
\noindent
Rád bych touto cestou srdečně poděkoval vedoucímu bakalářské práce Ing.~Markovi Beránkovi,~Ph.D.
za~čas věnovaný vedením této práce, za~příjemnou spolupráci a~za cenné poskytnuté rady a připomínky.

%%% první stránka
    %%%
%%%  PRVNÍ STRANA
%%%
%%%  * soubor obsahující první stranu bakalářské práce
%%%
%%%  ===========================================================================
\newpage
\pagestyle{plain}

%%% začátek stránkování
\setcounter{page}{6}

%%% nastavení speciálních okrajů
\newgeometry{textwidth=100mm, textheight=247mm, left=40.4mm, right=75mm}

%%% pozadí
\tikz[remember picture,overlay]
\node[opacity=1,inner sep=0pt] at (current page.center)
    {\includegraphics[width=\paperwidth,height=\paperheight]{\FIGDIR/side-banner}};

\begin{center}

    %%% logo školy
    \centerline{\mbox{\includegraphics[width=39.6mm]{\FIGDIR/uu-icon}}}

    \vfill

    %%% název práce v češtině
    \Large\textbf{Webové řešení na prodej vstupenek s~rezervací míst}

    \vspace{5mm}

    %%% název práce v angličtině
    \Large{Web-based ticketing solution with seat reservation}

    \vfill

    %%% logo školy
    \centerline{\mbox{\includegraphics[width=45mm]{\FIGDIR/uu-logo}}}

\end{center}
\restoregeometry

%%% abstrakt
    %%%
%%%  ABSTRAKT
%%%
%%%  * soubor obsahující abstrakt bakalářské práce
%%%
%%%  ===========================================================================
\newpage
\pagestyle{plain}

%%% abstrakt česky
\vbox to 0.5\vsize{
    \setlength\parindent{0mm}
    \setlength\parskip{5mm}

    %%% nadpis
    {\large\bfseries TODO: Abstrakt}

    %%% TODO: text abstrakt až po dopsání práce (lol)
    \noindent
    Tato bakalářská práce se zaměřuje na vývoj frontendové části webové aplikace pro prodej vstupenek s rezervací míst. Cílem práce je vyvinout prototyp aplikace, která potenciálnímu zákazníkovi umožní zobrazit mapu míst, vybrat si preferované místo, přidat vstupenky do nákupního košíku a následně vytvořit objednávku.

    Teoretická část práce se zaměřuje na obecnou problematiku prodeje vstupenek a moderního řešení pomocí platforem a služeb poskytujících online prodej vstupenek s možností rezervací míst. Tato část dále analyzuje trh současných poskytovatelů a definuje nejhlavnější technické problémy, které mohou při vývoji takového systému vzniknout a to výhradně z pohledu frontendu.

    Praktická část práce definuje rozsah implementované aplikace, podrobně popisuje hlavní funkce, komponenty, datové modely i některé nezbytné backendové části. Kapitola o návrhu uživatelského rozhraní popisuje principy, vzory a osvědčené postupy návrhu uživatelského rozhraní v kontextu vyvíjené aplikace. Kapitola o vývoji frontendové části poté podrobně popisuje technologie, nástroje a knihovny použité při vývoji aplikace.


    \vss}

%%% abstrakt anglicky
\nobreak\vbox to 0.49\vsize{
    \setlength\parindent{0mm}
    \setlength\parskip{5mm}

    %%% nadpis
    {\large\bfseries TODO: Abstract}

    %%% TODO: text anglicky
    \noindent
    This bachelor thesis deals with the implementation of a ticketing solution using seat reservations. The aim of
    the theoretical part of the thesis is to elaborate this issue and identify the main technical challenges and requirements for the implementation
    of such a system, especially from the frontend perspective. These challenges, such as intuitive user interface, availability of information
    real-time sales information, administration and management of seating plans or final booking and payment of orders, will be discussed in this thesis
    will be described in this thesis and options for their solution will be presented. The practical part will focus on the implementation of the part of the web application dealing with
    rendering of the interactive seating plan and its interaction with the customer. In this part the user interface design will be presented
    together with an introduction of the selected technologies to be implemented. Emphasis will be placed on the optimized design of the seating plan, the specification of
    of its data format, communcation with the API and overall documentation of the components used in the resulting application. The aim of this work is to describe
    the issues and challenges in implementing a web application addressing seat ticketing and present part of a real solution for such a
    application especially from the perspective of an interactive seat selection plan.

    Keywords: interactive seating plan, seat reservations, tickets, web applications, JavaScript, React, SVG
    \vss}


%%% obsah
    \newpage
    \openright
    \tableofcontents
%%% jednotlivé kapitoly
    \chapter{Úvod}

Používání tabulek a grafů v~odborném textu má některá společná
pravidla a~některá specifická. Tabulky a grafy neuvádíme přímo do
textu, ale umístíme je buď na samostatné stránky nebo na vyhrazené
místo v~horní nebo dolní části běžných stránek. \LaTeX\ se o~umístění
plovoucích grafů a tabulek postará automaticky. Každý graf a tabulku
očíslujeme a umístíme pod ně legendu. Legenda má popisovat obsah grafu
či tabulky tak podrobně, aby jim čtenář rozuměl bez důkladného
studování textu práce. Na každou tabulku a graf musí být v~textu odkaz
pomocí jejich čísla. Na příslušném místě textu pak shrneme ty
nejdůležitější závěry, které lze z~tabulky či grafu učinit.

\section{Podklady pro práci}

Lorem ipsum dolor sit amet, consectetuer adipiscing elit. Sed


%%% TODO: seznam použité literatury
%%% Reference se hledají v souboru priklady_literatury.bib. Aby se
%%% vytvořil seznam literatury, je třeba ocitovat alespoň jednu
%%% referenci, zkompilovat tento soubor latexem, pak bibtexem a znovu
%%% latexem. Tím se vytvoří seznam použitých referencí
%%% (BcPrace.bbl) a vloží se do práce na místě, kde se nachází příkaz
%%% \bibliography, tedy sem.
%    \bibliography{priklady_literatury}
%%% seznam obrázků
%    \listoffigures
%%% seznam tabulek
%    \listoftables

%%% TODO
%%% Použité zkratky v bakalářské práci, existují-li, včetně jejich vysvětlení.
%    \chapter*{Seznam použitých zkratek}
%    \addcontentsline{toc}{chapter}{Seznam použitých zkratek}


%%% TODO
%%% Přílohy k bakalářské práci, existují-li (různé dodatky jako výpisy programů,
%%% diagramy apod.). Každá příloha musí být alespoň jednou odkazována z vlastního
%%% textu práce. Přílohy se číslují.
%    \chapter*{Přílohy}
%    \addcontentsline{toc}{chapter}{Přílohy}


\end{document}
